\chapter{Maryam Mirzakhani --- Fields Medal 2014}\label{ch2014maryammirzakhani:maryam-mirzakhani-fields-medal-2014-a-comprehensive-report-on-her-work-and-impact}

\begin{center}\rule{0.5\linewidth}{0.5pt}\end{center}

\textbf{Prepared:} July 2026 \textbf{Subject:} Maryam Mirzakhani (May
12, 1977 -- July 14, 2017, Tehran, Iran) \textbf{Position:} Professor of
Mathematics, Stanford University (2008--2017) \textbf{Doctoral Advisor:}
Curtis T. McMullen (Fields Medalist, 1998) \textbf{Fields Medal
Citation:} \emph{``For outstanding contributions to the dynamics and
geometry of Riemann surfaces and their moduli spaces.''}

\begin{center}\rule{0.5\linewidth}{0.5pt}\end{center}

\section{Table of Contents}\label{ch2014maryammirzakhani:table-of-contents}

\begin{enumerate}
\def\labelenumi{\arabic{enumi}.}
\tightlist
\item
 Introduction
\item
 Biographical Background and Education
\item
 Riemann Surfaces and Moduli Spaces: Historical Context
\item
 The Counting Problem for Simple Closed Geodesics
\item
 Weil-Petersson Volumes of Moduli Spaces
\item
 The Generalized McShane Identity
\item
 Mirzakhani's Volume Formula
\item
 Intersection Theory on Moduli Spaces
\item
 Growth of Simple Closed Geodesics
\item
 Kontsevich-Witten Theory and Mirzakhani's Proof
\item
 Ergodic Theory of the Earthquake Flow
\item
 Teichmüller Dynamics
\item
 The Magic Wand Theorem: Orbit Closures for SL(2,R)
\item
 Invariant Measures on Moduli Space
\item
 Large Genus Asymptotics
\item
 Random Hyperbolic Surfaces
\item
 Broader Impact and Legacy
\item
 Awards and Recognition
\item
 Selected Major Publications
\item
 Conclusion
\item
 References and Further Reading
\end{enumerate}

\begin{center}\rule{0.5\linewidth}{0.5pt}\end{center}

\section{1. Introduction}\label{ch2014maryammirzakhani:introduction}

Maryam Mirzakhani was awarded the Fields Medal in 2014 at the age of 37,
becoming the first woman and the first Iranian citizen to receive the
award. She was recognized for her significant contributions to the dynamics
and geometry of Riemann surfaces and their moduli spaces --- work that
bridged hyperbolic geometry, complex analysis, topology, and dynamical
systems with notable creativity and technical power.

Mirzakhani's work is characterized by its originality and its ability to
connect seemingly disparate fields. Her PhD thesis solved a decades-old
problem about the growth rate of simple closed geodesics on hyperbolic
surfaces by finding an unexpected relationship to the volumes of moduli
spaces. She then derived a recursive formula for these volumes that not
only solved the counting problem but also provided a new proof of the
Kontsevich-Witten formula --- a deep result in intersection theory on
moduli spaces.

In the second half of her career, Mirzakhani turned to the dynamics of
moduli spaces. In a monumental collaboration with Alex Eskin (and later
Amir Mohammadi), she proved that orbit closures of the SL(2,R) action on
moduli space are always complex-algebraic subvarieties --- a result of
such power that it became known as the ``Magic Wand Theorem.'' This
work, which ran over 200 pages, has been described as ``the beginning of
a new era'' in the field.

Mirzakhani's career was tragically cut short by breast cancer, which
claimed her life on July 14, 2017, at the age of 40. In her brief but
productive career, she transformed the study of Riemann surfaces and
moduli spaces, opening new frontiers that continue to be explored by
mathematicians around the world.

This report provides a detailed account of Mirzakhani's important
contributions, the mathematical context in which they arose, the
techniques she developed, and their lasting impact on mathematics.

\begin{center}\rule{0.5\linewidth}{0.5pt}\end{center}

\section{2. Biographical Background and
Education}\label{ch2014maryammirzakhani:biographical-background-and-education}

\subsection{2.1 Early Life in Tehran}\label{ch2014maryammirzakhani:early-life-in-tehran}

Maryam Mirzakhani was born on May 12, 1977, in Tehran, Iran. She was the
second of three children. Her father, Ahmad Mirzakhani, was an
electrical engineer, and her mother, Zahra Kazemi, was a homemaker. From
an early age, Mirzakhani showed exceptional intellectual curiosity. She
later recalled that her passion for mathematics grew out of a love for
storytelling: ``I would make up stories in my head about solving
problems. The process of finding a solution became like a narrative.''

\subsection{2.2 Olympiad Success}\label{ch2014maryammirzakhani:olympiad-success}

Mirzakhani attended the Farzanegan School in Tehran, a selective school
for girls operated by the National Organization for Development of
Exceptional Talents. In 1994, she became the first Iranian woman to
compete in the International Mathematical Olympiad (IMO), winning a gold
medal with a score of 41 out of 42. The following year, in 1995, she won
a second gold medal with a perfect score of 42 out of 42. This
achievement made her an inspiration to young women across Iran and the
Middle East.

\subsection{2.3 University Education}\label{ch2014maryammirzakhani:university-education}

Mirzakhani enrolled at the Sharif University of Technology in Tehran,
where she earned a Bachelor of Science degree in Mathematics in 1999.
During her undergraduate studies, she published her first research
papers, including work on graph theory and combinatorial designs.

In 1999, Mirzakhani moved to the United States to pursue graduate
studies at Harvard University. Her doctoral advisor was Curtis T.
McMullen, himself a Fields Medalist (1998). At Harvard, she was known
for her determination and relentless questioning. She often took her
class notes in Persian, translating on the fly as she learned advanced
mathematics in a new language.

\subsection{2.4 Doctoral Thesis}\label{ch2014maryammirzakhani:doctoral-thesis}

Mirzakhani's PhD thesis, ``Simple Geodesics on Hyperbolic Surfaces and
the Volume of the Moduli Space of Curves'' (2004), was a landmark
achievement. In it, she solved the problem of counting simple closed
geodesics on hyperbolic surfaces by relating it to the volumes of moduli
spaces. The thesis produced three papers published in the top
mathematics journals: Inventiones Mathematicae, the Journal of the
American Mathematical Society, and the Annals of Mathematics.

\subsection{2.5 Academic Career}\label{ch2014maryammirzakhani:academic-career}

After completing her PhD, Mirzakhani held a Clay Mathematics Institute
Research Fellowship and an assistant professorship at Princeton
University (2004--2008). In 2008, she joined the faculty at Stanford
University as a full professor of mathematics --- a position she held
until her death in 2017.

\subsection{2.6 Personal Life and
Passing}\label{ch2014maryammirzakhani:personal-life-and-passing}

In 2008, Mirzakhani married Jan Vondrák, a Czech-American mathematician
and computer scientist. They had one daughter, Anahita. Mirzakhani was
diagnosed with breast cancer in 2013. Despite undergoing treatment, the
cancer metastasized, and she died on July 14, 2017, at Stanford
Hospital. She was 40 years old.

\begin{center}\rule{0.5\linewidth}{0.5pt}\end{center}

\section{3. Riemann Surfaces and Moduli Spaces: Historical
Context}\label{ch2014maryammirzakhani:riemann-surfaces-and-moduli-spaces-historical-context}

\subsection{3.1 Riemann Surfaces}\label{ch2014maryammirzakhani:riemann-surfaces}

A Riemann surface is a one-dimensional complex manifold --- a surface
that locally looks like the complex plane and on which one can do
complex analysis. Despite being ``one-dimensional'' in the complex
sense, a Riemann surface is two-dimensional as a real manifold, making
it a surface in the ordinary sense.

Every Riemann surface of genus \(g \ge 2\) admits a hyperbolic metric: a
Riemannian metric of constant curvature \(-1\). This is the famous
Uniformization Theorem, which states that any Riemann surface of genus
\(g \ge 2\) can be realized as a quotient of the hyperbolic plane
\(\mathbb{H}^2\) by a discrete group of isometries.

\subsection{3.2 Moduli Spaces}\label{ch2014maryammirzakhani:moduli-spaces}

The moduli space \(\mathcal{M}_g\) is the space of all Riemann surfaces
of genus \(g\), up to biholomorphic equivalence. Understanding the
geometry and topology of moduli spaces is one of the central problems in
modern mathematics, with connections to algebraic geometry, topology,
number theory, and theoretical physics.

The moduli space \(\mathcal{M}_g\) is not compact, but it can be
compactified (the Deligne-Mumford compactification
\(\overline{\mathcal{M}}_g\)) by adding surfaces with nodes (simple
singularities). The study of intersection numbers on
\(\overline{\mathcal{M}}_g\) is closely related to string theory and
quantum gravity through the work of Kontsevich and Witten.

\subsection{3.3 Teichmüller Space}\label{ch2014maryammirzakhani:teichmuxfcller-space}

Before moduli space, there is Teichmüller space \(\mathcal{T}_g\), which
is a covering space of \(\mathcal{M}_g\). Teichmüller space parametrizes
Riemann surfaces with a marking (a choice of a basis for the fundamental
group), and it is homeomorphic to a ball of dimension \(6g-6\). The
mapping class group acts on Teichmüller space, and the quotient is
\(\mathcal{M}_g\).

\begin{center}\rule{0.5\linewidth}{0.5pt}\end{center}

\section{4. The Counting Problem for Simple Closed
Geodesics}\label{ch2014maryammirzakhani:the-counting-problem-for-simple-closed-geodesics}

\subsection{4.1 Geodesics on Hyperbolic
Surfaces}\label{ch2014maryammirzakhani:geodesics-on-hyperbolic-surfaces}

On a hyperbolic surface, a geodesic is a curve that is locally the
shortest path between any two of its points. A closed geodesic is a
geodesic that forms a closed loop. A closed geodesic is \textbf{simple}
if it does not intersect itself.

The study of geodesics on hyperbolic surfaces has a long history. The
``prime number theorem for geodesics'' (Margulis, 1960s) states that the
number of closed geodesics of length at most \(L\) grows exponentially:

\[
\pi(L) \sim \frac{e^{hL}}{hL}
\]

where \(h\) is the topological entropy.

\subsection{4.2 The Simple Geodesic
Problem}\label{ch2014maryammirzakhani:the-simple-geodesic-problem}

The counting problem for \textbf{simple} closed geodesics was much more
difficult. Because simple geodesics cannot intersect themselves, they
are more constrained than general geodesics. Intuitively, one might
expect the number of simple closed geodesics to grow polynomially rather
than exponentially.

This problem was considered one of the important open questions in
hyperbolic geometry. As Benson Farb of the University of Chicago put it,
simple closed geodesics are ``the key object to unlocking the structure
and geometry of the whole surface.''

\subsection{4.3 Mirzakhani's Solution}\label{ch2014maryammirzakhani:mirzakhanis-solution}

In her PhD thesis, Mirzakhani proved that the number of simple closed
geodesics of length at most \(L\) on a hyperbolic surface of genus \(g\)
is asymptotic to:

\[
c_g \cdot L^{6g-6}
\]

where \(c_g\) is a constant depending on the hyperbolic structure. This
polynomial growth contrasts sharply with the exponential growth of
general closed geodesics and confirmed a long-standing conjecture.

The exponent \(6g-6\) is the dimension of the moduli space
\(\mathcal{M}_g\), and this is no coincidence: Mirzakhani's proof
revealed a deep connection between counting geodesics and computing
volumes of moduli spaces.

\begin{center}\rule{0.5\linewidth}{0.5pt}\end{center}

\section{5. Weil-Petersson Volumes of Moduli
Spaces}\label{ch2014maryammirzakhani:weil-petersson-volumes-of-moduli-spaces}

\subsection{5.1 The Weil-Petersson
Metric}\label{ch2014maryammirzakhani:the-weil-petersson-metric}

The moduli space \(\mathcal{M}_g\) carries a natural Riemannian metric
called the Weil-Petersson metric. This metric is not complete (the
moduli space can be completed by adding surfaces with nodes), but it has
a finite volume form. The volume of \(\mathcal{M}_g\) with respect to
this metric is a fundamental invariant that has been studied for
decades.

\subsection{5.2 The Volume Problem}\label{ch2014maryammirzakhani:the-volume-problem}

Computing the Weil-Petersson volume of \(\mathcal{M}_g\) is a classical
problem. For small genus, the volumes were known: -
\(\mathcal{M}_{0,3}\): volume 1 (a point) - \(\mathcal{M}_{1,1}\):
volume \(\pi^2/6\) - \(\mathcal{M}_{2}\): volume \(\pi^4/270\)

But for general genus, the problem seemed intractable.

\subsection{5.3 Bordered Riemann
Surfaces}\label{ch2014maryammirzakhani:bordered-riemann-surfaces}

Mirzakhani's key insight was to consider not just moduli spaces of
closed surfaces, but moduli spaces of \textbf{bordered} Riemann surfaces
--- surfaces with geodesic boundary components. This generalization
allowed her to use an inductive approach: by varying the lengths of the
boundary components, she could relate the volumes for different genera
and numbers of boundaries.

\begin{center}\rule{0.5\linewidth}{0.5pt}\end{center}

\section{6. The Generalized McShane
Identity}\label{ch2014maryammirzakhani:the-generalized-mcshane-identity}

\subsection{6.1 The Original McShane
Identity}\label{ch2014maryammirzakhani:the-original-mcshane-identity}

Greg McShane had proved a notable identity for punctured hyperbolic
surfaces: the sum over all simple closed geodesics of a certain function
of their lengths equals a constant. Specifically, for a one-punctured
torus:

\[
\sum_{\gamma} \frac{1}{1 + e^{\ell(\gamma)}} = \frac{1}{2}
\]

where the sum is over all simple closed geodesics \(\gamma\).

\subsection{6.2 Mirzakhani's
Generalization}\label{ch2014maryammirzakhani:mirzakhanis-generalization}

Mirzakhani generalized McShane's identity to bordered hyperbolic
surfaces with geodesic boundary components. For a surface with boundary
components of lengths \(x_1, \ldots, x_n\), she proved:

\[
\sum_{\{\gamma_1, \gamma_2\} \in \mathcal{F}_i} D(\ell_{\gamma_1}, \ell_{\gamma_2}, x_i) + \sum_{j \neq i} \sum_{\gamma \in \mathcal{F}_{i,j}} R(\ell_{\gamma}, x_i, x_j) = x_i
\]

where \(D\) and \(R\) are explicit elementary functions.

\subsection{6.3 From Identity to Volume}\label{ch2014maryammirzakhani:from-identity-to-volume}

The generalized McShane identity is the crucial link between geodesics
and volumes. By integrating both sides over the moduli space, Mirzakhani
transformed the identity into a recursive formula for the Weil-Petersson
volumes.

\begin{center}\rule{0.5\linewidth}{0.5pt}\end{center}

\section{7. Mirzakhani's Volume
Formula}\label{ch2014maryammirzakhani:mirzakhanis-volume-formula}

\subsection{7.1 The Recursive Formula}\label{ch2014maryammirzakhani:the-recursive-formula}

Mirzakhani proved that the Weil-Petersson volumes
\(V_{g,n}(L_1, \ldots, L_n)\) of the moduli space of bordered Riemann
surfaces satisfy a recursive formula:

\[
V_{g,n}(L_1, \ldots, L_n) = \int_0^{\infty} \int_0^{\infty} \left( \cdots \right) \cdot V_{g-1, n+2}(\cdots) \, dx \, dy + \sum_{\substack{g_1+g_2=g \\ I \sqcup J = \{1,\ldots,n\}}} \int_0^{\infty} \int_0^{\infty} \left( \cdots \right) \cdot V_{g_1, |I|+1}(\cdots) \cdot V_{g_2, |J|+1}(\cdots) \, dx \, dy
\]

This formula expresses the volume in terms of volumes of simpler moduli
spaces (with smaller genus or fewer boundary components).

\subsection{7.2 Polynomial Nature}\label{ch2014maryammirzakhani:polynomial-nature}

A key consequence of the recursion is that \(V_{g,n}(L_1, \ldots, L_n)\)
is a polynomial in \(L_1^2, \ldots, L_n^2\) with rational coefficients.
The top degree term is:

\[
V_{g,n}(L_1, \ldots, L_n) = \frac{1}{2^{2g+n-3} (3g-3+n)!} \cdot (\text{polynomial}) + \text{lower order terms}
\]

\subsection{7.3 The Constant Volumes}\label{ch2014maryammirzakhani:the-constant-volumes}

For closed surfaces (\(n=0\)), the formula gives:

\[
V_g(0) = \frac{1}{2^{2g-3} (3g-3)!} \cdot \sum_{m=0}^{3g-3} \frac{B_{2m}}{m!} \cdot \cdots
\]

where \(B_{2m}\) are Bernoulli numbers. For small \(g\): -
\(V_2 = \pi^4/270\) - \(V_3 = 11\pi^6/37800\) -
\(V_4 = 191\pi^8/11907000\)

\begin{center}\rule{0.5\linewidth}{0.5pt}\end{center}

\section{8. Intersection Theory on Moduli
Spaces}\label{ch2014maryammirzakhani:intersection-theory-on-moduli-spaces}

\subsection{8.1 The Kontsevich-Witten
Formula}\label{ch2014maryammirzakhani:the-kontsevich-witten-formula}

In the early 1990s, Maxim Kontsevich proved a notable formula
(conjectured by Edward Witten) relating the intersection theory of
moduli spaces to the KdV hierarchy of integrable partial differential
equations. This result, for which Kontsevich was awarded the Fields
Medal in 1998, is one of the deepest connections between mathematics and
physics.

The Kontsevich-Witten formula expresses the generating function for
intersection numbers of tautological classes on
\(\overline{\mathcal{M}}_{g,n}\) as a solution to the KdV hierarchy.

\subsection{8.2 Mirzakhani's New Proof}\label{ch2014maryammirzakhani:mirzakhanis-new-proof}

Mirzakhani discovered that her volume formula provides a new proof of
the Kontsevich-Witten formula. The connection comes through the
relationship between Weil-Petersson volumes and intersection numbers:

\[
V_{g,n}(L_1, \ldots, L_n) = \sum_{d_1 + \cdots + d_n \le 3g-3+n} \frac{L_1^{2d_1} \cdots L_n^{2d_n}}{2^{d_1+\cdots+d_n} d_1! \cdots d_n!} \cdot \langle \tau_{d_1} \cdots \tau_{d_n} \rangle_g
\]

where \(\langle \tau_{d_1} \cdots \tau_{d_n} \rangle_g\) are the
intersection numbers on \(\overline{\mathcal{M}}_{g,n}\).

By plugging her recursive formula for volumes into this expansion,
Mirzakhani derived recursive relations for the intersection numbers that
are equivalent to the KdV equations. This gave a completely new,
geometric proof of the Kontsevich-Witten formula --- one that avoided
the matrix integrals and combinatorial methods used by Kontsevich.

\subsection{8.3 Impact}\label{ch2014maryammirzakhani:impact}

Mirzakhani's proof was immediately recognized as a major contribution.
It showed that the Kontsevich-Witten formula, far from being an isolated
result, is intimately connected to the hyperbolic geometry of Riemann
surfaces. The proof also provided new tools for computing intersection
numbers and studying the topology of moduli spaces.

\begin{center}\rule{0.5\linewidth}{0.5pt}\end{center}

\section{9. Growth of Simple Closed
Geodesics}\label{ch2014maryammirzakhani:growth-of-simple-closed-geodesics}

\subsection{9.1 The Asymptotic Formula}\label{ch2014maryammirzakhani:the-asymptotic-formula}

Using her volume formulas, Mirzakhani proved the asymptotic:

\[
N_{g,n}(L) \sim \frac{L^{6g-6+2n}}{(6g-6+2n) \cdot \text{vol}(\mathcal{M}_{g,n})}
\]

for the number of simple closed geodesics of length at most \(L\) on a
hyperbolic surface of genus \(g\) with \(n\) cusps.

\subsection{9.2 The Constant}\label{ch2014maryammirzakhani:the-constant}

The constant \(c_g\) in the asymptotic formula is:

\[
c_g = \frac{1}{(6g-6) \cdot V_g(0)}
\]

where \(V_g(0)\) is the Weil-Petersson volume of \(\mathcal{M}_g\).

\subsection{9.3 Generalization of the Three Geodesics
Theorem}\label{ch2014maryammirzakhani:generalization-of-the-three-geodesics-theorem}

This result generalizes the classical ``theorem of the three geodesics''
for the sphere, which states that on any surface of genus 0 (a sphere),
there are exactly three simple closed geodesics. For higher genus
surfaces, the number grows polynomially, and Mirzakhani's work gave the
precise growth rate.

\begin{center}\rule{0.5\linewidth}{0.5pt}\end{center}

\section{10. Kontsevich-Witten Theory and Mirzakhani's
Proof}\label{ch2014maryammirzakhani:kontsevich-witten-theory-and-mirzakhanis-proof}

\subsection{10.1 Background on Intersection
Numbers}\label{ch2014maryammirzakhani:background-on-intersection-numbers}

On the moduli space \(\overline{\mathcal{M}}_{g,n}\), there are natural
cohomology classes called \(\psi\)-classes. For each marked point \(i\),
the class \(\psi_i\) is the first Chern class of the line bundle whose
fiber is the cotangent space at the \(i\)-th marked point. The
intersection numbers:

\[
\langle \tau_{d_1} \cdots \tau_{d_n} \rangle_g = \int_{\overline{\mathcal{M}}_{g,n}} \psi_1^{d_1} \cdots \psi_n^{d_n}
\]

are fundamental invariants of moduli spaces.

\subsection{10.2 Witten's Conjecture}\label{ch2014maryammirzakhani:wittens-conjecture}

In 1990, Edward Witten conjectured that the generating function for
these intersection numbers satisfies the KdV hierarchy --- a system of
nonlinear PDEs that arises in integrable systems theory. Witten's
conjecture sparked an explosion of research connecting moduli spaces,
integrable systems, and string theory.

\subsection{10.3 Kontsevich's Proof}\label{ch2014maryammirzakhani:kontsevichs-proof}

Kontsevich proved Witten's conjecture in 1992 using matrix models and
combinatorial methods. His proof was a tour de force that introduced
powerful new techniques, but it relied on a detour through mathematical
physics rather than working directly with moduli spaces.

\subsection{10.4 Mirzakhani's Geometric
Proof}\label{ch2014maryammirzakhani:mirzakhanis-geometric-proof}

Mirzakhani's proof was substantially different and more geometric. She
showed that: 1. The Weil-Petersson volumes are polynomials in the
boundary lengths. 2. These polynomials satisfy a recursive formula
derived from the generalized McShane identity. 3. The coefficients of
these polynomials are precisely the intersection numbers. 4. The
recursion for volumes translates into the KdV equations for the
generating function.

This proof connected hyperbolic geometry to enumerative geometry in a
completely unexpected way and is considered one of the most beautiful
results in modern geometry.

\begin{center}\rule{0.5\linewidth}{0.5pt}\end{center}

\section{11. Ergodic Theory of the Earthquake
Flow}\label{ch2014maryammirzakhani:ergodic-theory-of-the-earthquake-flow}

\subsection{11.1 Earthquake Deformations}\label{ch2014maryammirzakhani:earthquake-deformations}

An earthquake on a hyperbolic surface is a deformation that cuts the
surface along a simple closed geodesic, twists one side relative to the
other, and re-glues. The earthquake flow is a dynamical system on
Teichmüller space that describes how surfaces change under such
deformations.

\subsection{11.2 Thurston's Earthquake
Theorem}\label{ch2014maryammirzakhani:thurstons-earthquake-theorem}

William Thurston proved that any two hyperbolic surfaces can be
connected by a unique earthquake path. This fundamental result shows
that earthquakes are a complete deformation mechanism for hyperbolic
surfaces.

\subsection{11.3 Ergodicity of the Earthquake
Flow}\label{ch2014maryammirzakhani:ergodicity-of-the-earthquake-flow}

In 2008, Mirzakhani proved that the earthquake flow on Teichmüller space
is ergodic --- that is, any invariant measurable set has either zero or
full measure. This proved a conjecture of Thurston and showed that the
earthquake flow mixes the space of hyperbolic surfaces in a strong
statistical sense.

The proof used the connection between earthquake deformations and the
Weil-Petersson geometry of moduli spaces, combining tools from
hyperbolic geometry, dynamical systems, and the theory of measured
laminations.

\begin{center}\rule{0.5\linewidth}{0.5pt}\end{center}

\section{12. Teichmüller Dynamics}\label{ch2014maryammirzakhani:teichmuxfcller-dynamics}

\subsection{12.1 The Teichmüller Geodesic
Flow}\label{ch2014maryammirzakhani:the-teichmuxfcller-geodesic-flow}

Teichmüller space carries a natural metric (the Teichmüller metric) and
a dynamical system called the Teichmüller geodesic flow. This flow
describes how the complex structure of a Riemann surface changes along
geodesics in Teichmüller space.

The study of Teichmüller dynamics is intimately connected to the study
of billiards in polygons: the trajectory of a billiard ball on a
polygonal table can be ``unfolded'' to a geodesic on a Riemann surface,
and ergodic properties of the billiard correspond to properties of the
Teichmüller flow.

\subsection{12.2 The SL(2,R) Action on Moduli
Space}\label{ch2014maryammirzakhani:the-sl2r-action-on-moduli-space}

A key tool in Teichmüller dynamics is the action of the group SL(2,R) on
the moduli space of Riemann surfaces with a holomorphic 1-form (the
Hodge bundle). This action, discovered by Kontsevich and Zorich,
provides a general framework for studying the dynamics of moduli
spaces.

\subsection{12.3 Mirzakhani's
Contributions}\label{ch2014maryammirzakhani:mirzakhanis-contributions}

Mirzakhani made several fundamental contributions to Teichmüller
dynamics: - New counting results for closed geodesics in moduli space
(with Eskin) - Analysis of the distribution of simple closed geodesics -
Applications of the dynamics to the geometry of moduli spaces

Her work in this area culminated in the Magic Wand Theorem, her most
famous result.

\begin{center}\rule{0.5\linewidth}{0.5pt}\end{center}

\section{13. The Magic Wand Theorem: Orbit Closures for
SL(2,R)}\label{ch2014maryammirzakhani:the-magic-wand-theorem-orbit-closures-for-sl2r}

\subsection{13.1 The Problem}\label{ch2014maryammirzakhani:the-problem}

The action of SL(2,R) on the moduli space of abelian differentials (the
Hodge bundle over \(\mathcal{M}_g\)) has orbits that can be of various
dimensions. Understanding the possible orbit closures is a fundamental
problem in dynamics, analogous to the problem of understanding invariant
submanifolds in symplectic geometry.

The naive expectation is that orbit closures are complicated subsets ---
perhaps fractal, like strange attractors in chaotic dynamical systems.
Proving that they have any regular structure at all seemed out of reach.

\subsection{13.2 The Magic Wand Theorem}\label{ch2014maryammirzakhani:the-magic-wand-theorem}

In 2013, Mirzakhani and Alex Eskin proved a result that stunned the
mathematical community: every orbit closure of the SL(2,R) action on
moduli space is a complex-algebraic subvariety.

More precisely, they proved:

\textbf{Theorem (Eskin-Mirzakhani):} Every SL(2,R)-invariant measure on
the moduli space of abelian differentials is a linear combination of
algebraic measures supported on complex-analytic subvarieties.

In a follow-up with Amir Mohammadi, they proved the stronger result:

\textbf{Theorem (Eskin-Mirzakhani-Mohammadi):} Every orbit closure of
the SL(2,R) action is an algebraic subvariety.

\subsection{13.3 Why ``Magic Wand''?}\label{ch2014maryammirzakhani:why-magic-wand}

The result was so unexpected and powerful that it became known as the
``Magic Wand Theorem.'' The name reflects the idea that Mirzakhani and
Eskin had waved a magic wand and revealed that the seemingly chaotic
dynamics of moduli spaces are governed by rigid algebraic structures.

\subsection{13.4 The Proof}\label{ch2014maryammirzakhani:the-proof}

The proof, which spans over 200 pages, is a monumental achievement
combining: - Techniques from ergodic theory and dynamics - Algebraic
geometry and Hodge theory - Homogeneous dynamics and Ratner's theorem -
Teichmüller theory and the geometry of moduli spaces

The core idea was to use a ``measure rigidity'' argument, showing that
any invariant measure must be algebraic. This required developing new
techniques for studying measures in infinite-dimensional settings.

\subsection{13.5 Impact}\label{ch2014maryammirzakhani:impact-1}

The Magic Wand Theorem is a prominent achievement in dynamical systems in the 21st century. It has opened up
entirely new directions in the study of moduli spaces and has been
applied to problems in billiards, interval exchange transformations, and
the geometry of flat surfaces.

\begin{center}\rule{0.5\linewidth}{0.5pt}\end{center}

\section{14. Invariant Measures on Moduli
Space}\label{ch2014maryammirzakhani:invariant-measures-on-moduli-space}

\subsection{14.1 The Classification
Problem}\label{ch2014maryammirzakhani:the-classification-problem}

Following the Magic Wand Theorem, Eskin and Mirzakhani undertook the
even more ambitious project of classifying all invariant measures for
the SL(2,R) action. This resulted in a second monumental paper,
``Invariant and Stationary Measures for the SL(2,R) Action on Moduli
Space,'' published posthumously in 2018.

\subsection{14.2 Stationary Measures}\label{ch2014maryammirzakhani:stationary-measures}

A stationary measure is a measure that is invariant on average under a
random walk on the group. Eskin and Mirzakhani proved that stationary
measures for the SL(2,R) action are also algebraic --- that is, they are
linear combinations of algebraic measures supported on complex-analytic
subvarieties.

\subsection{14.3 Consequences}\label{ch2014maryammirzakhani:consequences}

The classification of invariant and stationary measures has significant
consequences: - It implies that the action of SL(2,R) on moduli space is
``rigid'' in a strong sense. - It provides a complete description of all
possible ergodic behaviors. - It has applications to the study of
billiards, translation surfaces, and interval exchange transformations.

\begin{center}\rule{0.5\linewidth}{0.5pt}\end{center}

\section{15. Large Genus Asymptotics}\label{ch2014maryammirzakhani:large-genus-asymptotics}

\subsection{15.1 Asymptotic Volumes}\label{ch2014maryammirzakhani:asymptotic-volumes}

In her 2013 paper ``Growth of Weil-Petersson Volumes and Random
Hyperbolic Surfaces of Large Genus,'' Mirzakhani studied the behavior of
Weil-Petersson volumes as the genus grows large. She proved:

\[
V_g \sim \frac{1}{\sqrt{\pi}} \left(\frac{2}{\pi^2}\right)^{g} (2g)!
\]

as \(g \to \infty\). This asymptotic formula shows that the volumes grow
super-exponentially with genus.

\subsection{15.2 Random Hyperbolic
Surfaces}\label{ch2014maryammirzakhani:random-hyperbolic-surfaces}

The volume asymptotic allows one to study ``random'' hyperbolic surfaces
of large genus. Mirzakhani showed that a random hyperbolic surface
(chosen with respect to the Weil-Petersson volume form) has certain
typical properties: - The length of the shortest geodesic is roughly
\(2\log(2g)\). - The Cheeger constant is bounded away from zero. - The
spectral gap of the Laplacian is bounded below.

These results inaugurated the study of random hyperbolic surfaces, which
has become an active area of research.

\begin{center}\rule{0.5\linewidth}{0.5pt}\end{center}

\section{16. Random Hyperbolic
Surfaces}\label{ch2014maryammirzakhani:random-hyperbolic-surfaces-1}

\subsection{16.1 The Weil-Petersson
Model}\label{ch2014maryammirzakhani:the-weil-petersson-model}

The moduli space \(\mathcal{M}_g\) with the Weil-Petersson volume form
provides a natural probability distribution on the set of hyperbolic
surfaces of genus \(g\). This is analogous to the Gaussian Unitary
Ensemble (GUE) in random matrix theory.

\subsection{16.2 Mirzakhani's Results}\label{ch2014maryammirzakhani:mirzakhanis-results}

Mirzakhani proved several fundamental results about random hyperbolic
surfaces:

\begin{itemize}
\tightlist
\item
 \textbf{Systole:} The expected length of the shortest closed geodesic
 is asymptotic to \(2\log(2g)\).
\item
 \textbf{Spectral Gap:} Random surfaces have a spectral gap of at least
 \(3/16\) with high probability.
\item
 \textbf{Mixing:} The geodesic flow on a random surface is
 exponentially mixing.
\end{itemize}

\subsection{16.3 Impact}\label{ch2014maryammirzakhani:impact-2}

This work opened up the field of ``random hyperbolic geometry,'' which
uses probabilistic methods to study the typical behavior of hyperbolic
surfaces. It has connections to random matrix theory, number theory
(through the distribution of eigenvalues of the Laplacian), and
mathematical physics.

\begin{center}\rule{0.5\linewidth}{0.5pt}\end{center}

\section{17. Broader Impact and Legacy}\label{ch2014maryammirzakhani:broader-impact-and-legacy}

\subsection{17.1 Impact on Mathematics}\label{ch2014maryammirzakhani:impact-on-mathematics}

Mirzakhani's work transformed several areas of mathematics: -
\textbf{Hyperbolic geometry:} Her counting results for simple closed
geodesics resolved a central problem. - \textbf{Moduli spaces:} Her
volume formulas and their connection to intersection theory provided new
tools and insights. - \textbf{Dynamics:} The Magic Wand Theorem is one
of the important results in dynamical systems. -
\textbf{Teichmüller theory:} Her work on the earthquake flow and
Teichmüller dynamics redefined the field.

\subsection{17.2 Impact Beyond
Mathematics}\label{ch2014maryammirzakhani:impact-beyond-mathematics}

As the first woman to win the Fields Medal, Mirzakhani became a symbol
of possibility for women and girls in mathematics. Her story --- from
Tehran to the Fields Medal --- inspired countless young people around
the world to pursue mathematics.

\subsection{17.3 The Mirzakhani Legacy}\label{ch2014maryammirzakhani:the-mirzakhani-legacy}

Following her death, several initiatives were established in her memory:
- \textbf{The Maryam Mirzakhani New Frontiers Prize} (Breakthrough Prize
Foundation): Awarded annually to early-career women mathematicians. -
\textbf{May 12 Initiative:} Each year on her birthday, events around the
world celebrate women in mathematics. - \textbf{The Mirzakhani Society:}
An organization at Stanford supporting women in mathematics.

\subsection{17.4 A Personal Reflection}\label{ch2014maryammirzakhani:a-personal-reflection}

Mirzakhani's colleague Benson Farb captured her mathematical style:
``She had a kind of daring, a willingness to try anything, to follow any
idea wherever it led. She was not afraid of hard problems. She was not
afraid of anything.''

\begin{center}\rule{0.5\linewidth}{0.5pt}\end{center}

\section{18. Awards and Recognition}\label{ch2014maryammirzakhani:awards-and-recognition}

\textbf{2014:} - Fields Medal (International Mathematical Union) - Clay
Research Award

\textbf{2013:} - Ruth Lyttle Satter Prize in Mathematics (American
Mathematical Society) - Simons Investigator Award

\textbf{2010:} - Plenary Speaker, International Congress of
Mathematicians

\textbf{2009:} - Blumenthal Award for the Advancement of Research in
Pure Mathematics (American Mathematical Society)

\textbf{2006:} - Popular Science ``Brilliant 10'' Award

\textbf{2005:} - Clay Research Fellowship

\textbf{2004:} - Harvard University's Alice T. Schafer Prize

\textbf{1995:} - Gold Medal, International Mathematical Olympiad
(perfect score)

\textbf{1994:} - Gold Medal, International Mathematical Olympiad

\begin{center}\rule{0.5\linewidth}{0.5pt}\end{center}

\section{19. Selected Major
Publications}\label{ch2014maryammirzakhani:selected-major-publications}

\begin{enumerate}
\def\labelenumi{\arabic{enumi}.}
\item
 \textbf{``Simple geodesics and Weil-Petersson volumes of moduli spaces
 of bordered Riemann surfaces.''} \emph{Inventiones Mathematicae}
 167(1), 179--222, 2007. --- \emph{The foundational paper introducing
 the volume formula and generalized McShane identity.}
\item
 \textbf{``Weil-Petersson volumes and intersection theory on the moduli
 space of curves.''} \emph{Journal of the American Mathematical
 Society} 20(1), 1--23, 2007. --- \emph{Connection between
 Weil-Petersson volumes and Kontsevich-Witten intersection numbers.}
\item
 \textbf{``Growth of the number of simple closed geodesics on
 hyperbolic surfaces.''} \emph{Annals of Mathematics} 168(1), 97--125,
 2008. --- \emph{Polynomial growth of simple closed geodesics.}
\item
 \textbf{``Ergodic theory of the earthquake flow.''}
 \emph{International Mathematics Research Notices} 2008(9), 2008. ---
 \emph{Proof of ergodicity of the earthquake flow (Thurston's
 conjecture).}
\item
 \textbf{``Growth of Weil-Petersson volumes and random hyperbolic
 surfaces of large genus.''} \emph{Journal of Differential Geometry}
 94(2), 267--300, 2013. --- \emph{Large genus asymptotics and random
 surfaces.}
\item
 \textbf{``Isolation, equidistribution, and orbit closures for the
 SL(2,R) action on moduli space''} (with A. Eskin and A. Mohammadi).
 \emph{Annals of Mathematics} 182(2), 673--721, 2015. --- \emph{The
 Magic Wand Theorem.}
\item
 \textbf{``Invariant and stationary measures for the SL(2,R) action on
 moduli space''} (with A. Eskin). \emph{Publications Mathématiques de
 l'IHÉS} 127, 95--324, 2018. --- \emph{Classification of invariant
 measures.}
\item
 \textbf{``Counting closed geodesics in Moduli space''} (with A.
 Eskin). \emph{Journal of Modern Dynamics} 5(1), 71--105, 2011. ---
 \emph{Asymptotic counting of Teichmüller geodesics.}
\item
 \textbf{``The boundary of an affine invariant submanifold''} (with A.
 Wright). \emph{Inventiones Mathematicae} 209(3), 927--984, 2017. ---
 \emph{Structure of boundaries of orbit closures.}
\item
 \textbf{``Towards large genus asymptotics of intersection numbers on
 moduli spaces of curves''} (with P. Zograf). \emph{Geometric and
 Functional Analysis} 25(4), 1258--1289, 2015. --- \emph{Asymptotic
 intersection numbers.}
\end{enumerate}

\begin{center}\rule{0.5\linewidth}{0.5pt}\end{center}

\section{20. Conclusion}\label{ch2014maryammirzakhani:conclusion}

Maryam Mirzakhani was one of the most productive and original
mathematicians of her generation. In a career that spanned barely 15
years, she transformed our understanding of Riemann surfaces, moduli
spaces, and the dynamics of Teichmüller space. Her contributions
include:

\begin{itemize}
\tightlist
\item
 Solving the counting problem for simple closed geodesics on hyperbolic
 surfaces
\item
 Deriving a recursive formula for Weil-Petersson volumes
\item
 Providing a new geometric proof of the Kontsevich-Witten formula
\item
 Proving ergodicity of the earthquake flow
\item
 Proving the Magic Wand Theorem on orbit closures of the SL(2,R) action
\item
 Classifying invariant measures on moduli space
\item
 Founding the study of random hyperbolic surfaces of large genus
\end{itemize}

Each of these achievements would be sufficient for a distinguished
career; together, they represent one of the most notable bodies of
work in modern mathematics.

The Fields Medal citation --- \emph{``for outstanding contributions to
the dynamics and geometry of Riemann surfaces and their moduli spaces''}
--- captures the breadth and depth of her work. Her insights have
integrated methods from hyperbolic geometry, complex analysis, algebraic
geometry, topology, and dynamics, opening new frontiers that
mathematicians will explore for generations to come.

Mirzakhani's tragic death at the age of 40 cut short a career of
notable promise. But the work she accomplished in her brief life
--- and the inspiration she provided as the first woman to win the
Fields Medal --- ensures that her legacy will endure.

\begin{center}\rule{0.5\linewidth}{0.5pt}\end{center}

\section{21. References and Further
Reading}\label{ch2014maryammirzakhani:references-and-further-reading}

\subsection{Official Sources}\label{ch2014maryammirzakhani:official-sources}

\begin{itemize}
\tightlist
\item
 IMU Fields Medal 2014 Announcement:
 https://www.mathunion.org/imu-awards/fields-medal/fields-medals-2014
\item
 Fields Medal Citation for Maryam Mirzakhani:
 https://www.mathunion.org/fileadmin/IMU/Prizes/Fields/2014/news\_release\_mirzakhani.pdf
\end{itemize}

\subsection{Biographical and Overview
Articles}\label{ch2014maryammirzakhani:biographical-and-overview-articles}

\begin{itemize}
\tightlist
\item
 Klarreich, E. ``A Tenacious Explorer of Abstract Surfaces.''
 \emph{Quanta Magazine}, August 12, 2014.
\item
 Stanford News. ``Stanford's Maryam Mirzakhani Wins Fields Medal.''
 August 12, 2014.
\item
 Roberts, S. ``Maryam Mirzakhani, First Woman to Win Fields Medal, Dies
 at 40.'' \emph{New York Times}, July 16, 2017.
\item
 Jackson, A. ``The Work of Maryam Mirzakhani.'' \emph{Notices of the
 AMS}, 2014.
\end{itemize}

\subsection{Key Technical References}\label{ch2014maryammirzakhani:key-technical-references}

\begin{itemize}
\tightlist
\item
 Mirzakhani, M. ``Simple geodesics and Weil-Petersson volumes of moduli
 spaces of bordered Riemann surfaces.'' \emph{Inventiones Mathematicae}
 167 (2007), 179--222.
\item
 Mirzakhani, M. ``Weil-Petersson volumes and intersection theory on the
 moduli space of curves.'' \emph{JAMS} 20 (2007), 1--23.
\item
 Mirzakhani, M. ``Growth of the number of simple closed geodesics on
 hyperbolic surfaces.'' \emph{Annals of Mathematics} 168 (2008),
 97--125.
\item
 Mirzakhani, M. ``Ergodic theory of the earthquake flow.'' \emph{IMRN}
 2008.
\item
 Eskin, A. and Mirzakhani, M. ``Invariant and stationary measures for
 the SL(2,R) action on moduli space.'' \emph{Publications Mathématiques
 de l'IHÉS} 127 (2018), 95--324.
\item
 Eskin, A., Mirzakhani, M., and Mohammadi, A. ``Isolation,
 equidistribution, and orbit closures for the SL(2,R) action on moduli
 space.'' \emph{Annals of Mathematics} 182 (2015), 673--721.
\end{itemize}
